Our guarantees rest on one structural fact---the null we plant is exact---and on the
supermartingale property of betting. The proofs are short, and their brevity is the
point: the design does the work, so we give every proof in full. Throughout,
$\F_{t-1}$ denotes the verifier's filtration of Section~\ref{sec:filtration},
generated by everything opened through step $t-1$; all bets and mixture weights are
$\F_{t-1}$-measurable (predictable).

A word on what changed in this version, since it is the technical heart of the paper.
An earlier draft defined the certified quantity as the \emph{marginal} sign
probability $\Pr[Z_i^{(s)}=1]$ and proved coverage from a supermartingale argument
that in fact requires the stronger conditional statement
$\Efrak[Z_i^{(s)}\mid\F_{i-1}] \ge p_0$ at every step. The two coincide when signs
are independent and identically distributed across pairs, and they come apart when
they are not---which is exactly our situation, since every canary passes through one
jointly trained and unlearned model. Section~\ref{sec:exp-dependence} measures the
gap and finds it large. We therefore restate the certificate against the quantity the
protocol actually controls, and we obtain a guarantee that needs no independence,
no exchangeability across pairs, and no homogeneity across templates or repetition
strata.

\subsection{The exact null}

\begin{theorem}[Exact conditional null]\label{thm:null}
Suppose the unlearned model $M_u$ is conditionally independent of the coin vector
$b=(b_1,\dots,b_m)$ given the pair contents $C=\{(c_i^0,c_i^1)\}$; exact unlearning
and retraining-from-scratch both satisfy this. Then for every score $s\in\F$, every
step $t$, and the pair $j=\pi(t)$ opened at that step,
\[
Z_j^{(s)} \mid \F_{t-1} \;\sim\; \Bern(\tfrac12) \quad\text{exactly,}
\]
for any model, any score, and any set of query prompts. In particular the sign
sequence is a sequence of exact conditional fair coins \emph{whether or not the signs
are independent across pairs}.
\end{theorem}

\begin{proof}
Enlarge $\F_{t-1}$ to $\G_{t-1} = \F_{t-1} \vee \sigma\bigl(\pi(t),\,
\{c_j^0, c_j^1\}\bigr)$, which reveals the position and the \emph{unordered} pair
opened at step $t$ but not its coin; this is legitimate because the manifest is
committed leaf-by-leaf (Section~\ref{sec:filtration}), so $b_j$ is opened only at
step $t$. Given $\G_{t-1}$ the two numbers $s(M_u, x)$ and $s(M_u, y)$ attached to
the unordered pair $\{x,y\}=\{c_j^0,c_j^1\}$ are determined. The coin $b_j$ is
independent of the other coins and of $C$ by construction, and independent of $M_u$
given $C$ by hypothesis; nothing opened before step $t$ involves $b_j$. Hence
$b_j \mid \G_{t-1} \sim \Bern(\tfrac12)$. Since $b_j$ decides which of the two
determined numbers is the in-twin's,
\[
D_j^{(s)} \;=\; s(M_u, c_j^{\text{in}}) - s(M_u, c_j^{\text{ghost}})
\;=\; \pm\bigl|s(M_u,x) - s(M_u,y)\bigr|
\]
with the two signs equally likely, so $D_j^{(s)} \overset{d}{=} -D_j^{(s)}$ given
$\G_{t-1}$. Therefore $\Pr[D_j^{(s)}>0\mid\G_{t-1}] = \Pr[D_j^{(s)}<0\mid\G_{t-1}]$,
and the committed tie-break coin splits any atom at zero equally, giving
$Z_j^{(s)}\mid\G_{t-1}\sim\Bern(\tfrac12)$. Averaging over $\G_{t-1}\supseteq\F_{t-1}$
gives the claim. No property of the model, the score, the query set, or the
dependence between pairs was used beyond measurability.
\end{proof}

Theorem~\ref{thm:null} is the reason VOUCH needs neither shadow models nor a
retrained reference: the null distribution is not estimated or assumed but generated
by our own coins, and it is exact at finite $m$. It is worth being explicit about why
the dependence objection does not bite here. The signs $Z_1^{(s)},\dots,Z_m^{(s)}$
are indeed dependent---they share one model---but the proof never asks for
independence. It asks only that the coin governing the pair opened at step $t$ still
be fair given everything opened before, and under exact unlearning it is, because the
model carries no information about it. The revocation arm therefore inherits an exact
null with no further conditions. The certificate arm, whose null is the composite
``some score still leaks,'' is where dependence has to be confronted, and that is the
next subsection.

\subsection{The certified quantity and the certificate}\label{sec:certquantity}

Once $M_u$ and the manifest exist, the sign vector $z^{(s)} = (z_1^{(s)},\dots,
z_m^{(s)})$ of the committed cohort is fixed; the only randomness left to the audit
is the reveal order and the tie coins. We therefore certify the advantage the
declared score actually achieves against that cohort.

\begin{definition}[Realised cohort advantage]\label{def:delta}
For $s\in\F$ the \emph{realised cohort advantage} is
\begin{equation}\label{eq:delta}
\Delta^{(s)} \;=\; \frac{2}{m}\sum_{i=1}^{m} z_i^{(s)} \;-\; 1
\;=\; 2\,\Pr\nolimits_{I}\bigl[z_I^{(s)}=1\bigr] - 1 ,
\end{equation}
the advantage of the rule ``guess that the twin scoring higher under $s$ is the
in-twin'' when the pair index $I$ is drawn uniformly from the committed cohort. Two
coarser quantities sit above it: the \emph{mechanism-level} advantage
$\Delta^{(s)}_C = \Efrak_b[\Delta^{(s)} \mid C, \xi]$, averaging over the inclusion
coins with the pair contents $C$ and the training randomness $\xi$ held fixed; and the
\emph{super-population} advantage $\bar\Delta^{(s)} = \Efrak[\Delta^{(s)}]$, averaging
over $C$ and $\xi$ as well.
\end{definition}

$\Delta^{(s)}$ is the quantity an auditor of \emph{this} model wants bounded, and it
is the one the protocol controls exactly; Definition~\ref{def:cert} is stated for it.
Proposition~\ref{prop:transfer} climbs one level, from $\Delta^{(s)}$ to
$\Delta^{(s)}_C$, at a stated cost in tolerance. We do not climb the second level to
$\bar\Delta^{(s)}$: that would need concentration over the draw of corpora, which we
neither prove nor assume. Where the experiments report a ``super-population'' verdict
they mean the fixed-boundary process of the earlier design, whose null is the
predictable conditional mean; we keep both columns throughout so that the reader can
see which claims survive at which tolerance.

\begin{definition}[$(\eps,\alpha)$-forgetting certificate]\label{def:cert}
Fix a tolerance $\eps \in (0,1)$, an error budget $\alpha \in (0,1)$, and the boundary
$p_0 = \tfrac12 + \tfrac\eps2$. An $(\eps,\alpha)$-forgetting certificate for $M_u$
over the class $\F$ is a data-dependent stopping decision that asserts
$\sup_{s\in\F}|\Delta^{(s)}| < \eps$, issued by a procedure with the guarantee that,
whenever in truth $\sup_{s\in\F}|\Delta^{(s)}| \ge \eps$, the probability that a
certificate is \emph{ever} issued---at any stopping time, under any monitoring
rule---is at most $\alpha$.
\end{definition}

The two-sidedness in Definition~\ref{def:cert} is deliberate: $\Delta^{(s)} \le -\eps$,
the over-forgetting regime in which the in-twin scores conspicuously \emph{below} its
never-seen partner, is membership leakage exactly as much as $\Delta^{(s)} \ge \eps$
is, and Section~\ref{sec:exp} shows that gradient-ascent unlearning lands there in
practice.

\begin{theorem}[Certificate coverage under arbitrary dependence]\label{thm:cert}
Reveal the cohort in a uniformly random order $\pi$ drawn by the verifier
\emph{independently of the realised sign vectors} $\{z^{(s)}\}_{s\in\F}$ and committed
before any query, and for each score $s$ and direction run the wealth process of
Eq.~\eqref{eq:wealthwor} with stakes confined to
$\lambda_t^{s,+} \in [0,\, c/(1-p_0(t)))$ and
$\lambda_t^{s,-} \in [0,\, c/\bar p_0(t))$ for some fixed $c < 1$, against the
predictable without-replacement boundary
\begin{equation}\label{eq:worboundary}
p_0(t) \;=\; \frac{m\,p_0 \;-\; \sum_{u<t} Z^{(s)}_{\pi(u)}}{m-t+1},
\end{equation}
declaring the null refuted outright if ever $p_0(t) > 1$. Issue the certificate only
when \emph{every} per-score, per-direction process exceeds $1/\alpha$. Then for any
class $\F$ and any stopping rule,
\[
\Pr\bigl[\text{a false }(\eps,\alpha)\text{-certificate is ever issued}\bigr] \le \alpha
\qquad\text{whenever } \sup_{s\in\F}|\Delta^{(s)}| \ge \eps .
\]
The guarantee is exact and finite-sample, and it holds \emph{conditionally on the
cohort's sign vectors}: no independence across pairs, no exchangeability across
pairs, and no homogeneity across templates or repetition strata is assumed. No
multiplicity correction over $\F$ is required. The one thing that \emph{is} assumed
is the independence of $\pi$ from the signs, and Remark~\ref{rem:order} explains why
it is not free.
\end{theorem}

\begin{proof}
Suppose the certification claim is false, so some score $s^\ast$ and direction violate
the tolerance; say $\Delta^{(s^\ast)} \ge \eps$, i.e.\ the cohort mean
$\mu = \tfrac1m\sum_i z_i^{(s^\ast)}$ satisfies $\mu \ge p_0$ (the case
$\Delta^{(s^\ast)} \le -\eps$ is symmetric, run against $1-p_0$). Condition on the
entire sign vector $z^{(s^\ast)}$; the remaining randomness is the uniform
permutation $\pi$, and sequential reveal of a uniform permutation is precisely simple
random sampling without replacement. Hence the conditional mean of the next sign is
the mean of the unrevealed remainder,
\[
\Efrak\bigl[Z^{(s^\ast)}_{\pi(t)} \,\big|\, \F_{t-1}\bigr]
= \frac{m\mu - \sum_{u<t} Z^{(s^\ast)}_{\pi(u)}}{m-t+1}
\;\ge\; \frac{m p_0 - \sum_{u<t} Z^{(s^\ast)}_{\pi(u)}}{m-t+1}
\;=\; p_0(t),
\]
using $\mu \ge p_0$. Each wealth factor therefore has conditional mean
\[
\Efrak\bigl[\,1 + \lambda_t\,(p_0(t) - Z^{(s^\ast)}_{\pi(t)}) \,\big|\, \F_{t-1}\bigr]
= 1 + \lambda_t\bigl(p_0(t) - \Efrak[Z^{(s^\ast)}_{\pi(t)} \mid \F_{t-1}]\bigr) \;\le\; 1,
\]
since $\lambda_t \ge 0$ is predictable. Nonnegativity is the other half and needs the
stake restriction: with $Z \in \{0,1\}$ the factor $1 + \lambda_t(p_0(t) - Z)$ is
minimised at $Z = 1$, where it equals $1 - \lambda_t(1 - p_0(t)) \ge 1 - c > 0$ by
$\lambda_t < c/(1-p_0(t))$; the mirror arm is identical with $\bar p_0(t)$ in place of
$1 - p_0(t)$. (This bound is why the implementation must derive the stake cap from the
\emph{same} clamped boundary the payoff is evaluated at; see
Appendix~\ref{app:repro}.) So the process is a nonnegative supermartingale started at
one, and Ville's inequality gives
$\Pr[\exists t: E_t^{s^\ast,+} \ge 1/\alpha] \le \alpha$. The refutation branch is
never taken under the null, because $p_0(t) \le \Efrak[Z_{\pi(t)}\mid\F_{t-1}] \le 1$.
Issuing the certificate requires \emph{every} process to exceed $1/\alpha$---in
particular the violating score's own---so a false certificate is issued with
conditional probability at most $\alpha$, uniformly over stopping rules; averaging
over the sign vectors preserves the bound. No union bound over $\F$ is needed because
the argument runs entirely through the single process attached to the violating
score.
\end{proof}

Two consequences are worth naming because they shape how the certificate should be
read. First, the guarantee spends no assumption on the \emph{data}: whatever
dependence the shared model induces among the signs, it is absorbed by conditioning on
the realised sign vector, and the only randomness the proof spends is the verifier's
own permutation. Second, and less comfortably, a cohort audited to exhaustion resolves
\emph{deterministically}: at $t=m$ the realised mean is known, and no inference
remains. The betting machinery therefore buys three things---early stopping at a
fraction of the query budget, validity under continuous monitoring and across
deletion waves, and the alarm arm, whose null is the genuinely stochastic one of
Theorem~\ref{thm:null}---rather than buying the cohort statement itself. We think
this is the honest way to describe what a canary audit can do, and
Section~\ref{sec:exp-baselines} quantifies each of the three.

\begin{remark}[The reveal order must be the verifier's own randomness]\label{rem:order}
Theorem~\ref{thm:cert} conditions on the sign vector and spends the permutation, so it
needs $\pi \perp \{z^{(s)}\}$. That is not a formality here. A provider
who can locate the cohort---and Section~\ref{sec:exp-detect} shows that is
cheap---could, if it also knew $\pi$, arrange which pairs are revealed early. The
released implementation therefore draws $\pi$ from one of three sources, and records
which on the certificate. A \emph{public randomness beacon} is the recommended one:
the verifier fetches a drand round published \emph{after} the manifest commitment and
seeds $\pi$ with $H(\text{root} \,\|\, \text{beacon randomness})$, so the order is
unpredictable to the provider before the round exists, independent of the signs, and
reconstructible afterwards by anyone holding the certificate---the audit's own
randomness becomes as checkable as its commitment. Failing network access, the
verifier draws from operating-system entropy, which satisfies the hypothesis but is
not third-party verifiable. The third source is the run seed, which is what produced
the results in this paper and which does \emph{not} satisfy the hypothesis, since that
seed also seeds cohort generation and training; nothing in the pipeline exploits the
coupling, but the condition is not met, and Appendix~\ref{app:repro} records exactly
where. Binding the beacon value to the commitment matters in both directions: drawing
the beacon first would let a provider grind the manifest against a known order.
\end{remark}

\begin{proposition}[Transfer to the coin-averaged advantage]\label{prop:transfer}
Fix the pair contents $C$ and the training randomness $\xi$, and write
$\Delta^{(s)}_C = \Efrak_b\bigl[\Delta^{(s)} \,\big|\, C, \xi\bigr]$ for the
advantage averaged over the inclusion coins alone---the advantage the unlearning
\emph{mechanism} yields on this corpus. Assume \emph{bounded coin influence}:
flipping a single coin $b_i$ changes at most $\kappa$ of the $m$ signs
$z_j^{(s)}$. Then for any $\alpha' \in (0,1)$, with probability at least
$1-\alpha'$,
\begin{equation}\label{eq:transfer}
\bigl|\Delta^{(s)} - \Delta^{(s)}_C\bigr| \;\le\;
\kappa\,\sqrt{\frac{2\log(2/\alpha')}{m}} .
\end{equation}
Consequently a cohort certificate at tolerance $\eps$ implies a mechanism-level
certificate at tolerance $\eps + \kappa\sqrt{2\log(2/\alpha')/m}$, at total error
level $\alpha + \alpha'$.
\end{proposition}

\begin{proof}
Write $f(b) = \tfrac2m\sum_j z_j^{(s)} - 1$ as a function of the coin vector, the
other randomness held fixed. Changing one coordinate changes $\sum_j z_j^{(s)}$ by at
most $\kappa$ by hypothesis, so $f$ has bounded differences $c_i = 2\kappa/m$.
McDiarmid's inequality gives
$\Pr[|f - \Efrak_b f| \ge u] \le 2\exp(-2u^2/\sum_i c_i^2) = 2\exp(-u^2 m/(2\kappa^2))$,
with the expectation over the coins and $(C,\xi)$ held fixed---which is exactly
$\Delta^{(s)}_C$; setting the right-hand side to $\alpha'$ yields
\eqref{eq:transfer}.
\end{proof}

Three qualifications travel with Proposition~\ref{prop:transfer} and we would rather
state them than let a reader over-read the bound. First, it bounds the distance to
$\Delta^{(s)}_C$, the \emph{coin-averaged} advantage on this corpus, not to
$\bar\Delta^{(s)} = \Efrak[\Delta^{(s)}]$ of Definition~\ref{def:delta}, which
additionally averages over the draw of $C$ and $\xi$; closing that second gap requires
a further concentration step over corpora that we do not have and do not claim.
Second, the bound is informative only when $\kappa = o(\sqrt m)$: at
$\kappa = \Theta(m)$ it is vacuous. Under exact unlearning $\kappa = 1$ holds
trivially, since a coin that leaves no trace in the model can only affect its own
pair's sign; away from it, an influential coin can move the model and hence many
signs, which is precisely the regime where the transfer is wanted. Third, the size of
the price is worth internalising: at $m=512$, $\alpha'=0.05$ and $\kappa=1$ the
inflation is $0.12$, which is why a mechanism-level claim at $\eps = 0.05$ needs a
cohort an order of magnitude larger than a cohort claim at the same tolerance---a gap
Section~\ref{sec:exp-bench} reports rather than elides.

\begin{remark}[Why the intuitive mixture fails]\label{rem:mixture}
It is tempting to run a \emph{single} certificate process on the adaptively reweighted
sign $\bar Z_t = \sum_s w_s Z_t^{(s)}$, and indeed this is valid for the revocation arm,
whose null makes every component an exact conditional fair coin. It is \emph{not}
valid for the certificate, whose composite null is ``some score still leaks.'' The
failure is easy to exhibit: take $\Delta^{(s_1)} = \eps$ exactly at the boundary and
$\Delta^{(s_2)} = \dots = 0$ clean. Then
$\Efrak[\bar Z_t \mid \F_{t-1}] = w_1 p_0 + (1 - w_1)/2 < p_0$ whenever $w_1 < 1$, so
each factor of the mixture's certificate process has conditional mean strictly greater
than one: the process is a \emph{sub}martingale under the null and crosses $1/\alpha$
with probability approaching one. Worse, the violation is generic rather than
adversarial, because discrimination-driven reweighting moves weight \emph{away} from
the one leaking score under partial unlearning. Against this construction---which is
our own earlier design, not a published method---we measured a false-certification
rate of $0.996$ against $0.007$ for the per-score minimum
(Table~\ref{tab:soundness}), which is why VOUCH requires consensus rather than an
average.
\end{remark}

\subsection{Power, confidence sequences, and composition}

\begin{theorem}[Power, fixed-boundary process]\label{thm:power}
Consider the \emph{super-population} arm, i.e.\ the process of
Eq.~\eqref{eq:wealthwor} run against the fixed boundary $p_0(t) \equiv p_0$ with signs
i.i.d.\ across pairs. If $|\Delta^{(s)}| < \eps$ for all $s$ and the bets have
sublinear regret (the mixture of Section~\ref{sec:engine} qualifies), then
$\tfrac1t \log E_t \to \KL(\Bern(p)\,\|\,\Bern(p_0))$ almost surely for each process,
and the expected certification time is
$\Efrak[\tau^*] = (1+o(1))\,\log(1/\alpha)/\min_{s,\pm}\KL$, the expectation being
taken under exact unlearning and the asymptotics in small $\eps$. Since
\begin{equation}\label{eq:klexp}
\KL\bigl(\Bern(\tfrac12)\,\big\|\,\Bern(\tfrac12 + \tfrac\eps2)\bigr)
\;=\; -\tfrac12\log\bigl(1-\eps^2\bigr)
\;=\; \tfrac{\eps^2}{2} + O(\eps^4),
\end{equation}
this is approximately $2\log(1/\alpha)/\eps^2$ pairs.
\end{theorem}

The restriction to the fixed boundary is not cosmetic. Under without-replacement reveal the log-wealth increments are
neither independent nor identically distributed, the boundary $p_0(t)$ moves, and the
admissible stake range moves with it---so ``the best constant bet'' is not well
defined, and neither the $O(\log t)$ ONS regret bound (which assumes a fixed convex
domain) nor the Krichevsky--Trofimov redundancy bound (which assumes a fixed $p_0$)
transfers as stated. Theorem~\ref{thm:power} should therefore be read as sizing the
cohort for a \emph{mechanism-level} claim. For the finite-cohort target we do not have
a matching analysis; its certification time is capped by $m$, since the claim resolves
deterministically once the cohort is exhausted, and we report its behaviour
empirically instead (Table~\ref{tab:power}, cohort column).

\begin{proof}
Fix one process, say the upper process for a score with true sign probability
$p < p_0$. For a constant bet $\lambda$, the log-wealth increments
$\log(1 + \lambda(p_0 - Z_t))$ are i.i.d.\ with mean
$f(\lambda) = p \log(1 - \lambda(1 - p_0)) + (1-p)\log(1 + \lambda p_0)$, so by the
strong law $\tfrac1t \log E_t \to f(\lambda)$ almost surely. The Kelly-optimal
constant bet is $\lambda^\ast = (p_0 - p)/(p_0(1 - p_0))$: substituting it turns the
two possible wealth factors into the Bernoulli likelihood ratios
$1 - \lambda^\ast(1 - p_0) = p/p_0$ and $1 + \lambda^\ast p_0 = (1-p)/(1-p_0)$,
so $f(\lambda^\ast) = p \log\tfrac{p}{p_0} + (1-p)\log\tfrac{1-p}{1-p_0}
= \KL(\Bern(p)\,\|\,\Bern(p_0))$, and no constant bet can do better because $f$ is
concave with this maximum. A betting strategy whose log-wealth regret against the best
constant bet is $o(t)$ therefore attains the same limit; the ONS expert has $O(\log t)$
regret, the KT plug-in \eqref{eq:kt} achieves $\KL - \tfrac{1}{2t}\log t$ directly by
the Krichevsky--Trofimov redundancy bound, and the uniform mixture over eight experts
trails its best expert by at most the constant $\log 8$, so all three satisfy the
hypothesis. For the crossing time, the process rejects when
$\log E_t \ge \log(1/\alpha)$; since $\tfrac1t\log E_t \to \KL$ almost surely, the
first-crossing time $\tau^\ast$ satisfies
$\Efrak[\tau^\ast] = (1 + o(1)) \log(1/\alpha)/\KL$ by the standard Wald argument, and
issuance waits for the \emph{slowest} process, whose drift is $\min_{s,\pm} \KL$.
Finally, under exact unlearning $p = \tfrac12$ for every score, and
\[
\KL\bigl(\Bern(\tfrac12)\,\big\|\,\Bern(p_0)\bigr)
= \tfrac12\log\tfrac{1/2}{p_0} + \tfrac12\log\tfrac{1/2}{1-p_0}
= -\tfrac12\log\bigl(4p_0(1-p_0)\bigr)
= -\tfrac12\log(1-\eps^2),
\]
whose expansion is \eqref{eq:klexp}, giving
$\Efrak[\tau^\ast] \approx 2\log(1/\alpha)/\eps^2$.
\end{proof}

Theorem~\ref{thm:power} tells the provider how large a canary cohort to plant. An
$\eps = 0.1$ certificate at $\alpha = 0.05$ needs on the order of $600$ pairs, an
$\eps = 0.05$ certificate roughly $2{,}400$, and an $\eps=0.02$ certificate roughly
$15{,}000$; because the drift equals the Kullback--Leibler divergence---the
information-theoretic optimum for testing $\Bern(\tfrac12)$ against $\Bern(p_0)$---no
valid procedure can do better by more than a constant factor. We note that an earlier
version of this theorem reported the expansion as $\eps^2/8$ and the sample size as
$8\log(1/\alpha)/\eps^2$; \eqref{eq:klexp} corrects it. The numerical columns of
Table~\ref{tab:power} were computed from the exact divergence and are unaffected.

\begin{theorem}[Finite-cohort certification time]\label{thm:cohortpower}
Fix a score $s$ and run the upper certificate process of
Theorem~\ref{thm:cert}: the without-replacement boundary
\eqref{eq:worboundary} with the rule that the null is refuted outright once
$p_0(t) > 1$. Let the committed cohort of size $m$ carry exactly $j$ signs
$z_i^{(s)} = 1$, and let $k_0 = \lceil m p_0 \rceil$ be the least number of
such signs compatible with the null $\Delta^{(s)} \ge \eps$. If
$j \le k_0 - 2$---the realised cohort sits at least two pairs below the
boundary---then the process refutes, and hence certifies at \emph{any}
$\alpha$, by the stopping time
\begin{equation}\label{eq:cohortworst}
\tau^\ast \;\le\; m - (k_0 - j) + 2
\qquad\text{for \emph{every} reveal order } \pi,
\end{equation}
and, under the uniformly random $\pi$ that Theorem~\ref{thm:cert} requires,
\begin{equation}\label{eq:cohortmean}
\Efrak[\tau^\ast] \;\le\; \frac{(m - k_0 + 1)(m+1)}{m - j + 1} \;+\; 1 ,
\end{equation}
with equality when the decision is reached by refutation. The lower arm is
the mirror statement with $z$ replaced by $1-z$ and $p_0$ by $1-p_0$; the
two-sided certificate over a class $\F$ is issued by the largest of the
$2|\F|$ individual bounds. Neither bound involves $\alpha$, the betting
strategy, or $|\F|$.
\end{theorem}

\begin{proof}
Refutation fires at the first $t$ with $p_0(t) > 1$, i.e.\ with
$m p_0 - \sum_{u<t} Z_{\pi(u)} > m - t + 1$. Writing $S_{t-1}$ for the
revealed count of ones and $W_{t-1} = (t-1) - S_{t-1}$ for the revealed
count of zeros, and using $S_{t-1} + (m - t + 1) = m - W_{t-1}$, the
condition is $m - W_{t-1} < m p_0$, that is $W_{t-1} > m(1 - p_0)$. Both
sides being compared to an integer, this is exactly
$W_{t-1} \ge m - k_0 + 1 =: r$, whether or not $m p_0$ is an integer.
Refutation therefore occurs one step after the $r$-th zero is revealed:
$\tau^\ast = 1 + \pi\text{-position of the } r\text{-th zero}$. The cohort
contains $m - j$ zeros, and $r \le m - j - 1$ exactly when $j \le k_0 - 2$,
so under that hypothesis the $r$-th zero is revealed at some position
$\le m - 1$ and $\tau^\ast \le m$.

For \eqref{eq:cohortworst}, the position of the $r$-th zero is largest when
every one precedes every zero, putting it at $j + r = j + m - k_0 + 1$;
adding the one step gives the bound, and the ordering that attains it shows
it is tight. For \eqref{eq:cohortmean}, let $M = m - j$ be the number of
zeros. In a uniformly random arrangement the $M$ zeros split the $m - M$
ones into $M+1$ exchangeable gaps, so the expected number of ones preceding
the $r$-th zero is $r(m - M)/(M + 1)$ and the expected position of that zero
is $r + r(m-M)/(M+1) = r(m+1)/(M+1)$. Adding one step gives
\eqref{eq:cohortmean}.

Finally, the argument never inspects the stake $\lambda_t$: refutation is a
statement about which cohort compositions remain feasible, so it holds for
every predictable betting strategy, and the certificate's actual stopping
time is the minimum of this time and the time its wealth crosses $1/\alpha$
by ordinary growth.
\end{proof}

Theorem~\ref{thm:cohortpower} is the finite-population counterpart that
Theorem~\ref{thm:power} could not supply. Under without-replacement reveal
the log-wealth increments are neither independent nor identically
distributed and the regret bounds behind Theorem~\ref{thm:power} do not
transfer, so we bound the cohort target's stopping time by a different
route: not by how fast wealth grows, but by when the null's remaining
feasible compositions run out. The two mechanisms run in parallel and the
process stops at whichever fires first, which is what makes the cohort
column of Table~\ref{tab:power} legitimately faster than the
Kullback--Leibler column rather than mysteriously so. Their regimes are
also easy to tell apart: exponential growth dominates when the cohort is
large relative to $\log(1/\alpha)/\KL$, and the feasibility bound binds when
it is not, which is exactly the tight-tolerance regime of
Section~\ref{sec:exp-tight}. At $m = 384$, $\eps = 0.2$ and a cohort at
chance ($j = m/2$), \eqref{eq:cohortmean} gives $308$ pairs against a
measured median of $165$: the wealth process usually wins, and the
feasibility bound is the guarantee that holds when it does not.

\begin{remark}[Per-score versus simultaneous confidence sequences]\label{rem:cs}
Each score's confidence sequence \eqref{eq:cs} covers that score's own advantage
$\Delta^{(s)}$ with uniform-in-time probability $1-\alpha$; the guarantee is
\emph{per score}. Two derived statements should be distinguished. The reported
$\max_{s\in\F} U_t^{(s)}$ \emph{is} a valid level-$\alpha$ anytime upper bound on
$\sup_{s\in\F}\Delta^{(s)}$ with no multiplicity correction, because the maximising
index of the population quantities is a fixed, non-random index once the cohort is
fixed, and that single score's own sequence is the only one the argument needs.
Simultaneous two-sided coverage of \emph{all} $|\F|$ advantages, on the other hand,
costs a union bound and holds at level $|\F|\alpha$; we report intervals at that
corrected level wherever we display the whole class.
\end{remark}

\begin{theorem}[Streaming composition, in both directions]\label{thm:streaming}
Let deletion waves $k = 1,\dots,K$ use cohorts with independent coins.
\emph{(a) Certificates compose by consensus.} Claiming that the whole history is clean
only when every wave has individually earned its certificate keeps the probability of a
false history-wide claim at or below $\alpha$ for any $K$; multiplying certificate
e-values instead tests the intersection null ``every wave is bad,'' whose rejection
certifies merely that \emph{some} wave is clean.
\emph{(b) Alarms compose by product.} Under the intersection null ``every wave was
exactly unlearned'' each per-wave revocation process is a supermartingale, so their
product is one too and crosses $1/\alpha$ with probability at most $\alpha$; the product
accumulates leakage spread so thinly across waves that no single wave reveals it.
\emph{(c)} Per-wave alarms with $\alpha$-spending $\alpha_k = \alpha 2^{-(k+1)}$ control
family-wise false revocation over an unbounded history.
\end{theorem}

\begin{proof}
\emph{(a)} Suppose the history-wide claim is false, so some wave $k^\ast$ truly has
$\sup_s |\Delta^{(s),(k^\ast)}| \ge \eps$. The all-pass rule claims only when every
wave has individually issued, so claiming requires wave $k^\ast$'s own per-score
minimum process to cross $1/\alpha$, an event of probability at most $\alpha$ by
Theorem~\ref{thm:cert}---independently of $K$ and of how the auditor schedules the
waves. For the negative half of the statement, note that a product
$\prod_k E^{(k)}$ of \emph{certificate} e-values is, by construction, an e-value for
the intersection null $\bigcap_k \{\text{wave } k \text{ bad}\}$; rejecting an
intersection establishes only the disjunction ``some wave is clean.'' In the
one-bad-wave scenario the clean waves' e-values grow without bound and drag the
product over $1/\alpha$ regardless of the bad wave. We measure the damage rather than
assert it: over $500$ ten-wave histories containing one genuinely bad wave, the
product rule issues a false history-wide certificate in \emph{every} history whose bad
wave truly violates the tolerance, against $0.000$ for the all-pass rule
(Table~\ref{tab:certprod}).

\emph{(b)} Under the intersection null, Theorem~\ref{thm:null} applies wave by wave,
so each $R^{(k)}$ is a nonnegative supermartingale with initial value one with respect
to the combined filtration, and cohorts across waves are independent by construction.
Conditional on $\F_{t-1}$, the increment of the product
$\prod_k R^{(k)}$ at any arrival is the increment of the single wave whose pair is
revealed, which has conditional mean at most one; the product is therefore itself a
nonnegative supermartingale started at one, for any adapted within-wave sampling times,
and Ville's inequality bounds its crossing probability by $\alpha$. Its power against
distributed leakage is additive in the exponent: if every wave carries advantage
$\Delta_0 > 0$, each factor drifts upward at rate
$\KL(\Bern(\tfrac12 + \tfrac{\Delta_0}2) \,\|\, \Bern(\tfrac12))$ per pair, so the
product crosses after roughly $\log(1/\alpha)/(m \cdot \KL)$ waves even when no single
wave is individually detectable.

\emph{(c)} Give wave $k$ its own alarm at level $\alpha_k = \alpha 2^{-(k+1)}$. By
Ville applied per wave, the probability that a truly clean wave $k$ ever alarms is at
most $\alpha_k$, and the union bound over an unbounded horizon costs
$\sum_{k \ge 1} \alpha 2^{-(k+1)} \le \alpha/2 \le \alpha$.
\end{proof}

\begin{proposition}[A magnitude-aware alarm]\label{prop:magnitude}
Under the null of Theorem~\ref{thm:null}, $\mathrm{sign}(D_t)$ is a fair coin even
conditional on $|D_t|$, so wealth factors of the form
$1 + \lambda_t\,\mathrm{sign}(D_t)\,g_t$ are an exact e-process for any predictable
stake $g_t \in [0,1]$ measurable with respect to the past and $|D_t|$. Taking $g_t$ to
be the predictable rank of $|D_t|$ concentrates bets on the large-effect pairs where
sparse memorisation lives, and halves detection time under $5\%$-sparse leakage
(Section~\ref{sec:exp-ablation}). Because it bets aggressively we report it as an
optional power enhancement and keep the sign-based arm as the default.
\end{proposition}

\begin{proof}
Symmetry of the conditional law of $D_t$ (Theorem~\ref{thm:null}) means that
conditioning on the magnitude $|D_t|$ leaves the sign exactly fair:
$\Efrak[\mathrm{sign}(D_t) \mid |D_t|, \F_{t-1}] = 0$, with ties randomised by the
committed coin. For any stake $g_t \in [0,1]$ that is measurable with respect to
$\F_{t-1} \vee \sigma(|D_t|)$---the predictable rank of $|D_t|$ among past magnitudes
qualifies---the tower property gives
$\Efrak[\mathrm{sign}(D_t)\, g_t \mid \F_{t-1}] = \Efrak[\,g_t\,
\Efrak[\mathrm{sign}(D_t) \mid |D_t|, \F_{t-1}] \mid \F_{t-1}] = 0$, so each factor
$1 + \lambda_t\, \mathrm{sign}(D_t)\, g_t$ with predictable $\lambda_t \in [0,1)$ has
conditional mean exactly one and stays positive. The product is a nonnegative
martingale with initial value one, hence an exact e-process, and Ville's inequality
applies as before. The power claim is empirical and quantified in
Section~\ref{sec:exp-ablation}.
\end{proof}

\begin{proposition}[One-directional bridge from certified unlearning]\label{prop:bridge}
If $\Ucal$ is $(\eps_u,\delta_u)$-certified, then
$|\bar\Delta^{(s)}| \le \tfrac{e^{\eps_u}-1}{e^{\eps_u}+1} + \delta_u$ for every $s$, so for
$\eps$ above this value VOUCH issues with probability $1 - \alpha - o(1)$. A certified
algorithm therefore passes VOUCH, as consistency demands, while VOUCH remains meaningful
at scales where $(\eps_u,\delta_u)$ certificates are unattainable.
\end{proposition}

\begin{proof}
The sign $Z^{(s)}$ is a randomised post-processing of the pipeline output: it is
computed from $M_u$ and the pair alone. $(\eps_u,\delta_u)$-certified unlearning makes
the law of $M_u$, under either value of the coin,
$(\eps_u,\delta_u)$-indistinguishable from the law of the retrained reference, under
which $Z^{(s)}$ is an exact fair coin by Theorem~\ref{thm:null}. Now view the rule
``guess that the twin with the larger score is the in-twin'' as a hypothesis test for
the coin $b_i$; its advantage is exactly
$\Pr[\text{correct}] - \Pr[\text{wrong}] = \bar\Delta^{(s)}$. The
trade-off region of \citet{kairouz2015composition} states that any test between two
$(\eps_u,\delta_u)$-indistinguishable distributions has type-I and type-II errors
satisfying $e^{\eps_u}\beta_1 + \beta_2 \ge 1 - \delta_u$ and
$\beta_1 + e^{\eps_u}\beta_2 \ge 1 - \delta_u$; adding the two constraints,
$\beta_1 + \beta_2 \ge 2(1 - \delta_u)/(1 + e^{\eps_u})$, so the advantage obeys
\[
|\bar\Delta^{(s)}| \;\le\; 1 - (\beta_1 + \beta_2)
\;\le\; 1 - \frac{2(1 - \delta_u)}{1 + e^{\eps_u}}
\;=\; \frac{e^{\eps_u} - 1}{e^{\eps_u} + 1} + \frac{2\delta_u}{1 + e^{\eps_u}}
\;\le\; \frac{e^{\eps_u} - 1}{e^{\eps_u} + 1} + \delta_u.
\]
When $\eps$ strictly exceeds this bound, every certificate process has strictly
positive drift by Theorem~\ref{thm:power} and crosses $1/\alpha$ with probability
tending to one as $m \to \infty$, while the revocation arm false-alarms with
probability at most $\alpha$; the issuance probability is therefore at least
$1 - \alpha - o(1)$.
\end{proof}

\begin{remark}[The bridge runs one way only]\label{rem:onedirection}
Proposition~\ref{prop:bridge} states that certified unlearning implies passing VOUCH.
The converse is false and we claim nothing like it: passing VOUCH does \emph{not}
imply any $(\eps_u,\delta_u)$ guarantee, does not bound the influence of any datum
outside the canary cohort, and does not constrain scores outside $\F$. The
proposition is a consistency check on the definition---an algorithm with a
parameter-space guarantee should not be flagged by a behavioural audit---and it
should not be read as a semantic equivalence in either direction.
\end{remark}

\paragraph{Scope, stated precisely.} A VOUCH certificate concerns the advantage
extractable through the declared class $\F$, on the committed canary cohort, against
the deployed model, under the honest-but-verifiable threat model of
Section~\ref{sec:setting}. It does not assert erasure in parameter space, which no
behavioural method can; it is a per-cohort statement, not a per-deletion-request one,
because per-request certificates would require per-user canaries; the transfer to the
super-population advantage costs the explicit inflation of
Proposition~\ref{prop:transfer}; and the extrapolation from planted canaries to
organic records is \emph{calibrated} through the dose--response strata and the
head-to-head comparisons of Section~\ref{sec:exp-descriptive} rather than proven. We
regard stating these limits precisely as part of the contribution, since a certificate
whose scope is overclaimed is worse than none.
